\section*{Chapter \ref{ch:m2:the-treasury-market} --- The Treasury Market}

\begin{solution}{exo:m2:the-treasury-market:1}
27 billion are left for competitive bids. 10 + 8 = 18 billion are filled below
4.110\%; the 10 billion at 4.110\% are needed for 9: stop-out 4.110\%,
allotment at the stop 90.00\%. Bid-to-cover $(34 + 3)/30 = 1.23$. No bidder
reaches the 35\% cap (10.5 billion).
\end{solution}

\begin{solution}{exo:m2:the-treasury-market:2}
$4.292 - 4.285 = +0.7$ basis points: a tail. $4.281 - 4.285 = -0.4$: a stop-through,
the strong result, since the Treasury sold at a lower yield (higher price)
than the market was paying a minute before.
\end{solution}

\begin{solution}{exo:m2:the-treasury-market:3}
4.125\% and 3.875\%: the eighths just below the stop-out yields. Priced at a
yield above its coupon, the note is issued slightly below par (99.409 for the
first, a ten-year).
\end{solution}

\begin{solution}{exo:m2:the-treasury-market:4}
The benchmark is the most liquid issue (tighter bid--ask, deeper books); it is
the one everybody hedges with and shorts, so it is often special in repo and
its holders finance it cheaply; and demand for the newest issue from
index-tracking and liquidity-seeking investors. The trade: buy the \omterm{def:m2:the-treasury-market:otr}{off-the-run},
short the \omterm{def:m2:the-treasury-market:otr}{on-the-run} in the same DV01, finance the long in repo and borrow the
short through a reverse repo, and wait for the new issue to become old at the
next auction. What goes wrong: the spread widens further in a flight to
liquidity (March 2020 did exactly this); the short becomes more special and the
financing costs more than the spread earns; leverage and margin force an exit
at the worst time.
\end{solution}

\begin{solution}{exo:m2:the-treasury-market:5}
$42\,000/26 = \text{USD}~1\,615$ million per dealer. The dealers' 38.3\% is what
the investors did not take: a high dealer share means the market's end
buyers were thin and the dealers must now distribute the notes, usually at a
price concession. The \omterm{def:m2:the-treasury-market:tail}{bid-to-cover ratio} counts every bid, including dealer
bids far above the stop-out that never had a chance to be filled.
\end{solution}

\begin{solution}{exo:m2:the-treasury-market:6}
The cap is 35\% of the offering: 14.7 billion. It can buy the remaining 5.3
billion in the secondary market, when-issued before the auction or after
issue, or from dealers who were awarded more than they wanted.
\end{solution}

\begin{solution}{exo:m2:the-treasury-market:7}
The stop-out moves from 4.198\% to 4.200\%, the tail from 1.8 to 2.0 basis
points: two tenths, not a whole basis point. The indirect bids are only part
of the book; as they move up, the dealer and direct bids that sat between
4.198 and 4.208 are reached first and fill most of the gap. The stop-out moves
by the full shift only if the whole book shifts.
\end{solution}

\begin{solution}{exo:m2:the-treasury-market:8}
In a \omterm{def:m2:the-treasury-market:auction}{uniform-price auction} a winning bidder pays the stop-out, not its bid. A
small bidder whose bid cannot move the stop-out has no reason to shade: it
should bid the yield at which it is indifferent, since bidding higher only
risks missing the award. A large bidder has a reason: a bid that sets the
stop-out moves the price it pays on its whole award, so it shades the part of
its demand that might be marginal. In a pay-your-bid auction everybody shades,
because every bid is its own price.
\end{solution}

\begin{solution}{pb:m2:the-treasury-market:1}
\textbf{1.} $42\,000 - 300 = \text{USD}~41\,700$ million.
\textbf{2.} Stop-out 4.198\%; tail $+1.8$ basis points.
\textbf{3.} Bid-to-cover 1.82; allotment at the stop 95.02\%.
\textbf{4.} Indirect 51.1\%, direct 10.6\%, dealers 38.3\%.
\textbf{5.} Coupon 4.125\%; price 99.409 per 100.
\textbf{6.} $42\,000/26 = 1\,615.4$ million; it bids 1\,615 million at three
yields, which meets the obligation if the yields are reasonably competitive.
\textbf{7.} The bids at 4.188\% and 4.195\% are below the stop-out and filled
in full; the bid at 4.205\% is above it and gets nothing: award 1\,000
million.
\textbf{8.} $807 \times 1\,000 = \text{USD}~806\,797$ per basis point.
\textbf{9.} It sold at 4.180\% and bought at 4.198\%: 1.8 basis points in its
favour, less 0.25 of costs, on 806\,797 per basis point: $+\text{USD}~1\,250\,536$.
\textbf{10.} A stop-through of 0.2 basis points: $808 \times 1\,000 \times (-0.2
- 0.25) = -\text{USD}~363\,758$.
\textbf{11.} At a tail of $+0.25$ basis points, the cost.
\textbf{12.} The end buyers bid for less than was offered at the market price,
so the Treasury had to sell deeper into the dealers' backstop bids; the
when-issued note, now the issued note, cheapens towards the stop-out as the
market absorbs the notes left with dealers.
\textbf{13.} To meet the pro-rata obligation, and because a bid a few tenths of
a basis point above the market is the price of underwriting: it wins if the
auction is weak, and then it is paid for the risk.
\textbf{14.} At 4.200\%, a tail of 2.0 basis points (\cref{exo:m2:the-treasury-market:7}).
\textbf{15.} It would have bid less and sold less when-issued, or bid lower
to be filled only in a weak auction. It cannot know the tail: the when-issued
level is the market's best estimate of the stop-out, and the dealer's P\&L
is a bet on the difference.
\textbf{16.} A uniform price removes most bidders' reason to shade, draws more
bidders (small bidders need not guess the others' bids), and makes every
winner pay the same price, the one the market cleared at. The Treasury has
used it for all its auctions since November 1998.
\textbf{17.} Some cannot bid directly or do not want to reveal their demand;
buying when-issued from a dealer fixes their price in advance, and the dealer
takes the auction risk for them.
\textbf{18.} In a weak auction it is filled on bids it would rather lose and
carries the notes at a loss until they are distributed. In return: \omterm{def:m2:the-treasury-market:pd}{primary
dealer} status, a counterparty relationship with the central bank, and client
flow.
\textbf{19.} Named result: \emph{the dealer's tail P\&L}: USD~806\,797 per basis
point of tail on its USD~1 billion award; USD~1\,250\,536 in this auction
(tail 1.8 basis points); break-even at a tail of 0.25 basis points.
\textbf{20.} The difference between what the auction paid and what the market
was paying at the deadline.
\end{solution}

\begin{solution}{iq:m2:the-treasury-market:1}
The stop-out yield minus the when-issued yield at the bid deadline. A large
positive tail means the auction had to reach up into higher yields than the
market was trading at to be filled: end demand was weak, dealers took more
than usual, and the market usually cheapens after the result.
\iqlookfor{the definition with its sign, and the link to dealer takedown.}
\end{solution}

\begin{solution}{iq:m2:the-treasury-market:2}
Liquidity (the benchmark is where trading, hedging and futures delivery
concentrate), specialness in repo (it is the most shorted issue, so owning
it lets you borrow cash cheaply), and demand from investors who need the
benchmark. Holders give up a little yield for these; the spread moves with
the premium on liquidity, and widens in stress.
\iqlookfor{liquidity and financing, not only ``it is newer''.}
\end{solution}

\begin{solution}{iq:m2:the-treasury-market:3}
In a pay-your-bid auction every winner pays its own bid, so bidders shade
towards the expected stop-out and must guess others' bids; winners regret
(winner's curse). In a \omterm{def:m2:the-treasury-market:auction}{uniform-price auction} winners pay the clearing yield,
so small bidders bid their true reservation yields; only bidders large enough
to set the price have a reason to shade.
\iqlookfor{who pays what, and which bidders still shade.}
\end{solution}

\begin{solution}{iq:m2:the-treasury-market:4}
After a weak retail-sales release, depth in the ten-year order books fell to a
fraction of normal; between 9:33 and 9:45 yields fell 16 basis points and
retraced with no news, while principal trading firms did more than half of
the volume. In order-book data: depth at the best levels and within a few
ticks over time, cancellation and order-to-trade ratios, the share of volume
by participant type, self-trades, aggressive versus passive volume by
participant, and the futures-cash lead--lag.
\iqlookfor{depth, not volume, as the variable, and a concrete list of
measurable quantities.}
\end{solution}

\begin{solution}{iq:m2:the-treasury-market:5}
Because everyone needed cash, and Treasuries are what can be sold: leveraged
investors unwinding, foreign holders raising dollars and funds meeting
redemptions sold them, and dealers could not absorb the flow within their
balance sheets. The safe asset was liquid, so it was sold first; the central
bank's purchases, over a trillion dollars in five weeks, restored function.
\iqlookfor{the distinction between safety and liquidity, and the
intermediation capacity of dealers.}
\end{solution}

\begin{solution}{iq:m2:the-treasury-market:6}
Entities: security (CUSIP, term, dated and maturity dates), announcement
(offering, dates), when-issued quotes (time-stamped, source), result (high
yield, allotment at high, bid-to-cover, bidder shares, coupon, price),
computed tail. Real-time path: subscribe to the result release, parse, look
up the when-issued quote at the deadline from a time-indexed store, compute,
publish. Tests: parse the published results of past auctions and compare with
their published statistics; a when-issued quote exactly at the deadline and
missing quotes; reopened issues (coupon already known, quoted in price);
bills (discount rates, not yields); clock and time-zone handling of the
deadline.
\iqlookfor{the deadline timestamp as the crux, and tests against published
history.}
\end{solution}
